\chapter{Cartes réseau et contournement du noyau}\label{ch:nw:network-cards-and-kernel-bypass}

Entre le câble et le \texttt{recv} d'une socket UDP, un datagramme est copié par la carte dans la mémoire de l'hôte, signalé par une
interruption, repris par le code réseau du noyau dans une interruption logicielle, rattaché à une socket, mis en file d'attente puis
recopié dans le tampon que le programme a remis à \texttt{recv}~; le programme, entre-temps, dormait et doit être réveillé. Sur
l'ordinateur portable où ce livre est rédigé, le datagramme médian passe moins d'une microseconde à entrer dans le noyau par la boucle locale, puis une quinzaine
de microsecondes de plus pour en sortir à destination d'un lecteur endormi. Rien de tout cela ne décide de quoi que ce soit qui intéresse un programme de trading.

Ce chapitre suit un paquet à travers la \omterm{def:nw:networking-for-trading:frame}{carte réseau} et le noyau, puis retire le noyau du chemin. Le
livre~13 (chapitre~13) a réglé l'hôte pour que la scrutation active et le contournement du noyau puissent jouer leur rôle~; ce chapitre examine la même
mécanique du côté de la carte~: anneaux de descripteurs, interruptions et leur regroupement, \omterm{def:nw:network-cards-and-kernel-bypass:rss}{aiguillage des flux} vers les files, piles
qui s'exécutent en espace utilisateur, et horodatages que la carte elle-même peut prendre.

\section{Le chemin d'un paquet à travers le noyau}

Une \omterm{def:nw:networking-for-trading:frame}{carte réseau} et son pilote partagent une mémoire. La carte écrit les paquets qui arrivent dans les tampons que le pilote lui a
confiés, et indique lesquels elle a remplis~; le pilote les prend, rend les tampons et remonte les paquets.

\begin{definition}[Anneau de descripteurs]\label{def:nw:network-cards-and-kernel-bypass:ring}
Un \emph{anneau de descripteurs}\index{anneau de descripteurs} est un tableau circulaire de descripteurs, chacun pointant vers un tampon de paquet dans la mémoire de l'hôte
et portant un drapeau de propriété. Côté réception, la carte remplit le prochain descripteur qui lui appartient avec un paquet arrivant (par accès direct
à la mémoire), le marque comme appartenant à l'hôte et avance sa tête~; l'hôte traite dans l'ordre les descripteurs qui lui appartiennent et les rend
à la carte. Lorsque le descripteur suivant appartient encore à l'hôte, la carte n'a nulle part où placer le paquet et le rejette.
\end{definition}

\begin{omfigure}
\begin{tikzpicture}[font=\footnotesize,b/.style={draw,rounded corners=2pt,minimum height=0.6cm,align=center,inner sep=3pt},>=Stealth]
\node[b,fill=gray!10] (w) at (0,0) {câble};
\node[b,fill=omDef!18] (n) at (2.0,0) {carte};
\node[b,fill=omDef!10] (r) at (4.4,0) {anneau de\\ descripteurs (DMA)};
\node[b,fill=omThm!12] (i) at (6.9,0.9) {interruption,\\ softirq};
\node[b,fill=omThm!12] (s) at (9.2,0.9) {file de\\ socket};
\node[b,fill=omThm!12] (c) at (11.4,0.9) {copie, réveil,\\ \texttt{recv}};
\node[b,fill=omMeth!20] (u) at (8.1,-1.1) {scrutation de l'anneau en espace utilisateur (contournement)};
\draw[->] (w) -- (n);
\draw[->] (n) -- (r);
\draw[->,omThm,thick] (r.north east) -- (i.west);
\draw[->,omThm,thick] (i) -- (s);
\draw[->,omThm,thick] (s) -- (c);
\draw[->,omMeth,thick] (r.south east) -- (u.west);
\node[font=\scriptsize,omThm] at (9.2,1.75) {chemin du noyau};
\node[font=\scriptsize,omMeth] at (8.1,-1.75) {chemin de contournement~: le programme possède l'anneau};
\end{tikzpicture}

\omcaption{Deux sorties d'un \omterm{def:nw:network-cards-and-kernel-bypass:ring}{anneau de descripteurs}. Le chemin du noyau (en haut) prend une interruption, traite le paquet dans une
interruption logicielle, le place dans une file de socket et le copie vers un programme qu'il faut réveiller. Une pile de contournement du noyau (en bas) projette l'anneau dans
le programme, qui le scrute et lit le paquet à l'endroit où la carte l'a écrit.}\label{fig:nw:network-cards-and-kernel-bypass:path}
\end{omfigure}

Sous Linux, la partie noyau de ce mécanisme est NAPI, le mécanisme d'événements de la pile réseau~: la carte lève une interruption,
le noyau planifie le traitement de l'anneau dans une interruption logicielle, et ce traitement se poursuit jusqu'à ce que l'anneau soit vide ou qu'un budget soit épuisé,
les nouvelles interruptions de cet anneau étant désactivées entre-temps. Sous charge, le noyau scrute donc déjà~; ce qu'il ne fait pas par défaut, c'est
épargner au programme la file de socket, la copie et le réveil.

\begin{omfigure}
\begin{tikzpicture}[font=\scriptsize,>=Stealth]
\foreach \k in {0,...,11} {
  \pgfmathsetmacro{\a}{90-30*\k}
  \pgfmathtruncatemacro{\host}{(\k>=2 && \k<=6) ? 1 : 0}
  \ifnum\host=1 \def\c{omThm!30}\else\def\c{omDef!15}\fi
  \draw[fill=\c] (\a+15:1.3) arc (\a+15:\a-15:1.3) -- (\a-15:2.2) arc (\a-15:\a+15:2.2) -- cycle;
}
\draw[->,very thick,omMeth] (0,0) -- (30:1.2);
\node[font=\scriptsize,omMeth] at (30:0.45) [above left] {queue};
\draw[->,very thick,omDef] (0,0) -- (-120:1.2);
\node[font=\scriptsize,omDef] at (-120:0.5) [left] {tête};
\node[anchor=west,align=left] at (2.8,0.9) {\tikz\fill[omThm!30] (0,0) rectangle (0.3,0.2);\ possédé par l'hôte~: rempli par la carte,\\
  \hphantom{\tikz\fill (0,0) rectangle (0.3,0.2);\ }en attente de traitement ou de restitution};
\node[anchor=west,align=left] at (2.8,-0.3) {\tikz\fill[omDef!15] (0,0) rectangle (0.3,0.2);\ possédé par la carte~: tampons vides};
\node[anchor=west,align=left] at (2.8,-1.2) {la carte remplit à la tête et avance dans le sens horaire,\\ l'hôte lit à la queue~;
  lorsque la tête atteint\\ un descripteur possédé par l'hôte, le paquet suivant est rejeté};
\end{tikzpicture}

\omcaption{Un anneau de réception de douze descripteurs. Cinq appartiennent à l'hôte (remplis et pas encore rendus), sept à la carte. La carte
écrit à la tête, l'hôte lit à la queue~; l'anneau est plein lorsque la tête rencontre un descripteur que l'hôte possède encore,
qu'il l'ait lu ou qu'il ne l'ait pas encore rendu.}\label{fig:nw:network-cards-and-kernel-bypass:ring}
\end{omfigure}

\begin{proposition}[Quand un anneau déborde]\label{prop:nw:network-cards-and-kernel-bypass:full}
Soit un anneau de $N$ descripteurs qui reçoit des paquets au débit $\lambda$ pendant que l'hôte n'en traite aucun pendant un temps $T$, à partir d'un anneau
où l'hôte détient $H$ descripteurs. L'anneau rejette $\bigl(\lambda T - (N - H)\bigr)^{+}$ paquets. Un
hôte qui rend les descripteurs par lots de taille $b$ en détient à tout instant jusqu'à $b - 1$ déjà traités, il doit donc
dimensionner l'anneau pour la plus longue pause augmentée de $b - 1$.
\end{proposition}

\begin{proof}
La carte dispose de $N - H$ descripteurs à remplir~; chaque arrivée en consomme un et aucun ne revient pendant la pause. Les descripteurs traités mais
pas encore rendus restent à l'hôte jusqu'à ce que le lot soit complet.
\end{proof}

\section{Interruptions, regroupement et répartition côté réception (RSS)}

Une interruption coûte à l'hôte un passage dans le noyau, le gestionnaire, et, pour un lecteur endormi, le réveil par l'ordonnanceur~: des microsecondes,
payées à chaque interruption. Les cartes attendent donc avant d'interrompre.

\begin{definition}[Regroupement des interruptions]\label{def:nw:network-cards-and-kernel-bypass:coalesce}
Le \emph{regroupement des interruptions}\index{regroupement des interruptions} (interrupt coalescing) est la pratique d'une carte qui retarde l'interruption de réception après
l'arrivée d'un paquet, jusqu'à ce qu'un délai donné soit écoulé ou qu'un nombre donné de paquets attendent, afin qu'une seule interruption serve plusieurs paquets.
Sous Linux, les deux limites se règlent avec \texttt{ethtool -C} sous les noms \texttt{rx-usecs} et \texttt{rx-frames}~; zéro
microseconde et une trame désactivent le regroupement.
\end{definition}

\begin{proposition}[Ce que le regroupement apporte et ce qu'il coûte]\label{prop:nw:network-cards-and-kernel-bypass:coalesce}
Soit des paquets arrivant selon un processus de Poisson de taux $\lambda$, la carte interrompant $u$ après le premier paquet en attente (la limite en trames
n'étant jamais atteinte), et soit $c_{\mathrm{irq}}$ le coût en CPU de chaque interruption et $c_{\mathrm{pkt}}$ celui de chaque paquet. Chaque interruption sert en
moyenne $1 + \lambda u$ paquets, l'hôte consomme une fraction
\[
  \frac{\lambda\, c_{\mathrm{irq}}}{1 + \lambda u} + \lambda\, c_{\mathrm{pkt}}
\]
d'un cœur, et le premier paquet de chaque lot attend $u$ pour l'interruption. Sans regroupement ($u = 0$), l'hôte sature un cœur
à $\lambda = 1/(c_{\mathrm{irq}} + c_{\mathrm{pkt}})$.
\end{proposition}

\begin{proof}
Un lot s'ouvre avec un paquet et dure $u$~; par la propriété d'absence de mémoire du processus de Poisson, $\lambda u$ paquets supplémentaires arrivent
en moyenne. Les interruptions surviennent donc au taux $\lambda/(1+\lambda u)$, chacune coûtant $c_{\mathrm{irq}}$, et chaque paquet coûte
$c_{\mathrm{pkt}}$.
\end{proof}

\begin{definition}[Répartition côté réception, aiguillage des flux]\label{def:nw:network-cards-and-kernel-bypass:rss}
La \emph{répartition côté réception}\index{répartition côté réception} (RSS) est la distribution, par une carte multi-files, des paquets arrivants sur plusieurs
anneaux de réception, selon un hachage de leurs adresses et de leurs ports, afin que différents cœurs puissent traiter différents flux. \emph{Aiguillage
des flux}\index{aiguillage des flux} est le placement par la carte de flux choisis (par correspondance exacte sur adresses et ports) sur un anneau choisi,
en passant outre le hachage.
\end{definition}

Une machine de trading utilise les deux à rebours de l'intention de RSS. RSS répartit la charge uniformément, ce que veut un serveur web~; un feed handler veut que
chaque paquet de la ligne A se trouve sur l'unique anneau que son cœur épinglé scrute, et rien d'autre. L'\omterm{def:nw:network-cards-and-kernel-bypass:rss}{aiguillage des flux} place les groupes de la plateforme sur les anneaux
des cœurs qui les traitent, et le reste du trafic de la machine ailleurs, et le cœur qui scrute un anneau doit se trouver sur le socket
auquel la carte est rattachée (livre~13, chapitre~4).

\section{Pilotes en mode scrutation et piles réseau en espace utilisateur}

\begin{definition}[Pilote en mode scrutation, pile réseau en espace utilisateur]\label{def:nw:network-cards-and-kernel-bypass:pmd}
Un \emph{pilote en mode scrutation}\index{pilote en mode scrutation} s'exécute dans le processus de l'application~: il projette les anneaux d'une carte dans l'espace utilisateur et
les lit et les remplit en boucle sur un cœur dédié, sans interruption ni appel système. Une \emph{pile réseau en espace
utilisateur}\index{pile réseau en espace utilisateur} implémente IP, UDP et TCP dans le processus de l'application, par-dessus un tel accès à la carte,
généralement derrière les appels de socket standard, de sorte qu'un programme non modifié s'y exécute.
\end{definition}

\begin{remark}[Les familles de contournement]\label{rem:nw:network-cards-and-kernel-bypass:families}
Trois familles sont en usage, et une firme en utilise souvent deux. L'accès brut aux anneaux par l'interface d'un fabricant (le \texttt{ef\_vi} d'AMD pour ses
cartes Solarflare) ou par un cadre de scrutation portable (DPDK) donne au programme des trames et rien d'autre~: il doit analyser son propre UDP et,
pour les ordres, apporter son propre TCP. Une pile en espace utilisateur derrière les appels de socket (OpenOnload, libre, qui intercepte les appels de socket
d'un binaire non modifié) conserve le code du programme et retire le noyau de son chemin. Et l'AF\_XDP propre à Linux déplace des trames de
l'anneau d'une carte vers une zone de la mémoire du programme au moyen d'un filtre dans le noyau, en mode copie ou en mode zero-copy, le noyau restant
chargé du périphérique. La \Cref{met:nw:network-cards-and-kernel-bypass:choose} indique lequel employer et où.
\end{remark}

\begin{method}[Choisir une voie de réception]\label{met:nw:network-cards-and-kernel-bypass:choose}
\begin{enumerate}
\item Données de marché sur le chemin critique (hot path)~: accès brut aux anneaux, scruté par un cœur épinglé sur le socket de la carte, le feed handler analysant lui-même l'UDP~; le
  flux est en multicast UDP, qui ne nécessite aucun état de protocole.
\item Envoi d'ordres~: une pile TCP en espace utilisateur, ou le TCP du fabricant de la carte implémenté en matériel (chapitre~7), car l'envoi d'ordres exige l'état de TCP et
  la retransmission, et écrire son propre TCP est un projet à part entière.
\item Tout le reste (drop copy, données de référence, supervision)~: le noyau, sur des cœurs et des anneaux qui lui sont propres, de sorte qu'il ne
  partage jamais une file avec le chemin critique.
\item Dans tous les cas, mesurer~: le gain dépend de la carte, du réglage de la machine et du trafic, et un scrutateur qui est un jour désordonnancé
  perd davantage qu'il n'a gagné.
\end{enumerate}
\end{method}

\omcode{code/firm/nicring/cpp/firm_nicring.hpp}{18}{37}{L'anneau de réception de \texttt{firm.nicring}~: la carte remplit les descripteurs qui lui appartiennent~;
la boucle de scrutation traite ceux qui lui appartiennent et les rend par lots.}\label{lst:nw:network-cards-and-kernel-bypass:ring}

\section{Réseau sans copie et côté émission}

\begin{definition}[Réseau sans copie]\label{def:nw:network-cards-and-kernel-bypass:zerocopy}
Le \emph{réseau sans copie}\index{réseau sans copie} est la remise d'un paquet au programme, ou son envoi par celui-ci, dans le tampon que la carte
lit ou écrit par accès direct à la mémoire (DMA), sans que le processeur recopie ses octets entre un tampon du noyau et celui du programme.
\end{definition}

Une copie d'un datagramme de 100 octets représente quelques nanosecondes de trafic mémoire~; ce que le \omterm{def:nw:network-cards-and-kernel-bypass:zerocopy}{réseau sans copie} économise du côté
réception tient moins à la copie qu'à la mécanique qui l'entoure (le tampon de socket, sa comptabilité, les lignes de cache d'un second
emplacement). Côté émission, la même structure d'anneaux fonctionne à l'envers~: le programme écrit un ordre dans un tampon d'émission, remplit un descripteur et écrit dans le registre de sonnette
de la carte, et la carte va chercher la trame par accès direct à la mémoire. Chacune de ces étapes traverse l'interconnexion du processeur vers
la carte (livre~13, chapitre~4). Certaines cartes permettent au programme de pousser les octets de la trame dans leur propre mémoire par des écritures
du processeur à la place~: le \texttt{ef\_vi} d'AMD documente une telle zone d'E/S programmées, qui «~évite la latence associée à une lecture DMA~»,
peut être remplie avant le chemin critique et mise à jour sur celui-ci, et, dans un mode dit cut-through, commence à émettre pendant que la trame est encore
en train de traverser le bus. L'émission la moins coûteuse est celle qui est préparée avant d'être nécessaire~: la trame de l'ordre
construite à l'avance, avec seulement le prix, la taille et l'identifiant à renseigner, l'idée même que le chapitre~7 reprend dans le matériel.

\section{Horodatages sur la carte}

Le noyau peut horodater un paquet à son entrée dans le noyau (\texttt{SO\_TIMESTAMPING}, horodatages de réception logiciels), et une carte
dotée de sa propre horloge peut l'horodater lorsqu'il atteint le port de la carte, un horodatage matériel au sens du livre~13, chapitre~5. La
différence est le temps passé par le paquet dans la carte et dans son anneau, ce qu'une firme qui règle son chemin de réception veut justement voir et
ce que les horodatages logiciels ne peuvent pas montrer. L'horloge de la carte n'est aussi bonne que sa synchronisation sur la
référence de temps de la firme, sujet du chapitre~4, et les horodatages de toute la baie sont comparés au chapitre~5.

\section{Tutoriel~: le chemin de réception du noyau, mesuré et modélisé}\label{tut:nw:network-cards-and-kernel-bypass:tut}

\begin{tutorial}
\textbf{Objectif.} Mesurer ce que coûte le chemin de réception du noyau sur l'ordinateur portable, modéliser ce que coûtent les interruptions et la scrutation sur une carte, et faire tourner
l'\omterm{def:nw:network-cards-and-kernel-bypass:ring}{anneau de descripteurs} dans trois langages. \textbf{État final~:} le \cref{tab:nw:network-cards-and-kernel-bypass:measured},
les \cref{fig:nw:network-cards-and-kernel-bypass:measured,fig:nw:network-cards-and-kernel-bypass:coalesce}, et le test de référence de l'anneau au vert
en Python, C++ et Rust.
\begin{tutsteps}
\item \textbf{Trois façons de lire.} \texttt{cpp/nw\_rx.hpp} envoie des datagrammes de 64 octets sur la boucle locale, chacun portant son
  heure d'envoi, et les lit par un \texttt{recvmsg} bloquant, par des lectures non bloquantes dans une boucle, et par \texttt{recvmmsg} par lots de 16, avec l'horodatage de réception logiciel
  du noyau dans chacun (le \cref{lst:nw:network-cards-and-kernel-bypass:rx}). \texttt{python bench\_rx.py} exécute 20\,000 datagrammes par mode,
  à 20 microsecondes d'intervalle, et écrit les quantiles.
\item \textbf{Découper le chemin.} L'horodatage du noyau découpe la latence de chaque datagramme en son entrée dans le noyau et sa sortie
  vers le programme.
\item \textbf{Modéliser la carte.} \texttt{nw\_nic.coalesce} et \texttt{poll} simulent les interruptions d'une carte avec regroupement et un
  cœur de scrutation, avec les coûts du modèle~; \texttt{python fig\_nic.py} écrit les données de la \cref{fig:nw:network-cards-and-kernel-bypass:coalesce}.
\item \textbf{L'anneau.} \texttt{firm.nicring} fait passer le même flux d'arrivées et de scrutations par quatre configurations d'anneau en Python,
  C++ et Rust, et vérifie les compteurs et les empreintes par rapport au test de référence commun.
\end{tutsteps}
\textbf{Pour aller plus loin.} Épingler le lecteur en attente active sur un cœur et l'expéditeur sur un autre (\texttt{taskset}), puis relire à nouveau sa queue
de distribution~; dans le test de référence de l'anneau, trouver le plus petit anneau qui ne perd aucun paquet lorsque la machine rend ses descripteurs par lots de 32.
\end{tutorial}

\omcode{code/networks/03-network-cards-and-kernel-bypass/cpp/nw_rx.hpp}{100}{116}{Les trois lectures~: une lecture bloquante, une lecture non bloquante
dans une boucle, et une lecture par lots~; chaque datagramme conserve ses heures d'envoi, du noyau et de l'application.}\label{lst:nw:network-cards-and-kernel-bypass:rx}

\begin{table}[ht]
\centering\small
\begin{tabular}{@{}>{\raggedright\arraybackslash}p{3.6cm}>{\centering\arraybackslash}p{2.6cm}>{\centering\arraybackslash}p{2.6cm}>{\centering\arraybackslash}p{2.1cm}>{\centering\arraybackslash}p{2.1cm}@{}}
\toprule
Méthode de lecture & Entrée dans le noyau, p50 & Du noyau au programme, p50 & Total, p50 & Total, p99 \\
\midrule
\texttt{recvmsg} bloquant & \qty{0.75}{\micro\second} & \qty{15.4}{\micro\second} & \qty{16.2}{\micro\second} & \qty{215}{\micro\second} \\
Attente active (busy polling) & \qty{0.62}{\micro\second} & \qty{0.92}{\micro\second} & \qty{1.5}{\micro\second} & \qty{3\,457}{\micro\second} \\
\texttt{recvmmsg} par lots & \qty{0.45}{\micro\second} & \qty{21.4}{\micro\second} & \qty{22.0}{\micro\second} & \qty{75}{\micro\second} \\
\bottomrule
\end{tabular}
\caption{Le chemin de réception du noyau sur la boucle locale, 20\,000 datagrammes de 64 octets par méthode, à 20 microsecondes d'intervalle (le
lecteur par lots reçoit des rafales de 16 envois). Mesuré sur un ordinateur portable (Intel Core Ultra 7 155H) sous WSL2, sans cœurs isolés~; g++~11,
\texttt{-O2}. Données~: \texttt{bench\_rx.py}, \texttt{measured\_rx.csv}.}\label{tab:nw:network-cards-and-kernel-bypass:measured}
\end{table}

\begin{omfigure}
\begin{tikzpicture}
\begin{axis}[width=11.5cm,height=5.4cm,ymode=log,xmin=0,xmax=1,ymin=0.7,ymax=6000,
  xlabel={quantile},ylabel={latence (\unit{\micro\second}, échelle logarithmique)},
  ytick={1,10,100,1000},yticklabels={1,10,100,1\,000},xtick={0,0.25,0.5,0.75,0.9,1},
  xticklabels={0,0.25,0.5,0.75,0.9,1},
  tick label style={font=\small},label style={font=\small},grid=major,grid style={gray!25},
  legend columns=3,legend style={font=\footnotesize,at={(0.5,-0.3)},anchor=north,draw=none,fill=none,
  /tikz/every even column/.append style={column sep=8pt}},table/col sep=comma]
\addplot[omThm,very thick,mark=*,mark size=1.3pt] table[x=q,y=blocking]{figdata/networks/03-network-cards-and-kernel-bypass/measured_rx_q.csv};
\addplot[omDef,very thick,mark=square*,mark size=1.3pt] table[x=q,y=busy]{figdata/networks/03-network-cards-and-kernel-bypass/measured_rx_q.csv};
\addplot[omMeth,very thick,mark=triangle*,mark size=1.6pt] table[x=q,y=batched]{figdata/networks/03-network-cards-and-kernel-bypass/measured_rx_q.csv};
\legend{lecture bloquante,attente active,lecture par lots}
\end{axis}
\end{tikzpicture}

\omcaption{Quantiles du temps écoulé entre \texttt{sendto} et le programme qui détient le datagramme, sur la boucle locale de l'ordinateur portable, selon la méthode de lecture.
L'attente active divise la médiane par dix~; sans cœur isolé, le thread de scrutation est parfois désordonnancé pendant des millisecondes, et chaque
datagramme qui le suit attend~: à partir du 95\up{e} quantile, c'est la pire des trois méthodes. Mesuré sur un ordinateur portable (Intel Core Ultra 7 155H) sous
WSL2, sans cœurs isolés. Données~: \texttt{bench\_rx.py}.}\label{fig:nw:network-cards-and-kernel-bypass:measured}
\end{omfigure}

La mesure fait apparaître les deux moitiés de l'argument de ce chapitre (\cref{tab:nw:network-cards-and-kernel-bypass:measured}). Faire entrer un
datagramme dans le noyau coûte moins d'une microseconde~; le faire sortir vers un programme endormi en coûte quinze à la médiane, et des centaines de
microsecondes au 99\up{e}~centile, parce que le programme doit être réveillé et ordonnancé. Un programme qui ne dort jamais s'en affranchit~: l'attente active
ramène la médiane à \qty{1.5}{\micro\second}. Et un programme qui ne dort jamais ne doit jamais être interrompu~: sur un ordinateur portable
dont les cœurs sont partagés avec tout le reste, le thread de scrutation est parfois désordonnancé, et les datagrammes qui arrivent entre-temps attendent
pendant des millisecondes. Le traitement par lots amortit l'appel système, mais non le réveil, et il attend que son lot soit complet. La boucle locale ne fait intervenir aucune \omterm{def:nw:networking-for-trading:frame}{carte réseau}~:
ces chiffres ne représentent donc que la part du noyau~; une carte réelle y ajoute son anneau, son interruption et son regroupement.

\omcode{code/networks/03-network-cards-and-kernel-bypass/python/nw_nic.py}{24}{45}{Le modèle de regroupement~: une interruption au $m$-ième
paquet en attente ou $u$ après le premier, puis le réveil, le coût de l'interruption et le coût par paquet.}\label{lst:nw:network-cards-and-kernel-bypass:model}

\begin{omfigure}
\begin{tikzpicture}
\begin{groupplot}[group style={group size=2 by 1,horizontal sep=1.8cm},width=6.6cm,height=5.0cm,xmin=0,xmax=66,
  xlabel={temporisation du regroupement (\unit{\micro\second})},tick label style={font=\small},label style={font=\small},
  grid=major,grid style={gray!25},table/col sep=comma,unbounded coords=jump]
\nextgroupplot[ymin=0,ymax=48,ylabel={latence moyenne (\unit{\micro\second})},xtick={0,8,16,32,64},
  legend columns=2,legend style={font=\footnotesize,at={(1.15,-0.32)},anchor=north,draw=none,fill=none,
  /tikz/every even column/.append style={column sep=8pt}}]
\addplot[omDef,very thick,mark=*,mark size=1.3pt] table[x=timer_us,y=lat_01]{figdata/networks/03-network-cards-and-kernel-bypass/coalesce.csv};
\addplot[omThm,very thick,mark=square*,mark size=1.3pt] table[x=timer_us,y=lat_05]{figdata/networks/03-network-cards-and-kernel-bypass/coalesce.csv};
\legend{100\,000 paquets par seconde,500\,000 paquets par seconde}
\nextgroupplot[ymin=0,ymax=0.7,ylabel={part d'un cœur},xtick={0,8,16,32,64}]
\addplot[omDef,very thick,mark=*,mark size=1.3pt] table[x=timer_us,y=cpu_01]{figdata/networks/03-network-cards-and-kernel-bypass/coalesce.csv};
\addplot[omThm,very thick,mark=square*,mark size=1.3pt] table[x=timer_us,y=cpu_05]{figdata/networks/03-network-cards-and-kernel-bypass/coalesce.csv};
\end{groupplot}
\end{tikzpicture}

\omcaption{Modèle~: \omterm{def:nw:network-cards-and-kernel-bypass:coalesce}{regroupement des interruptions} à deux débits de paquets. Des temporisations plus longues coûtent de la latence et économisent du processeur. À 500\,000 paquets par seconde, une
interruption par paquet sature le cœur (aucun point à zéro). Un cœur en scrutation, non tracé, maintient une latence voisine de
\qty{0.4}{\micro\second} aux deux débits et occupe le cœur entier. Coûts du modèle~: \qty{2}{\micro\second} par interruption,
\qty{3}{\micro\second} de réveil, \qty{0.3}{\micro\second} par paquet. Données~: \texttt{fig\_nic.py}.}\label{fig:nw:network-cards-and-kernel-bypass:coalesce}
\end{omfigure}

Le modèle (\cref{lst:nw:network-cards-and-kernel-bypass:model}) compte \qty{2}{\micro\second} de processeur par interruption,
\qty{3}{\micro\second} de réveil et \qty{0.3}{\micro\second} par paquet~: ce sont des valeurs de l'ordre du réveil de l'ordinateur portable, et non celles
d'une carte particulière. À 100\,000 paquets par seconde, une interruption par paquet donne une latence moyenne de \qty{5.6}{\micro\second} pour 23~\% d'un cœur~;
une temporisation de \qty{8}{\micro\second} donne \qty{11.7}{\micro\second} pour 14~\%~; la scrutation donne \qty{0.4}{\micro\second} pour le cœur entier.
À 500\,000 paquets par seconde, une interruption par paquet dépasse la capacité d'un cœur (le débit de saturation de la
\cref{prop:nw:network-cards-and-kernel-bypass:coalesce} est de 435\,000 par seconde), et le regroupement est ce qui maintient l'hôte en service. Un hôte de trading
résout ce compromis en ne le faisant pas~: le chemin critique (hot path) scrute sur des cœurs dédiés à cet usage, et tout le reste regroupe.

\section{Construction~: l'anneau de descripteurs}\label{bld:nw:network-cards-and-kernel-bypass:nicring}

\begin{build}
\bfield{Rôle} Un modèle de l'anneau de réception d'une carte, assez précis pour raisonner sur les pertes, le traitement par lots et les budgets de scrutation, dans
les trois langages des systèmes de la société.
\bfield{Interface} Python \texttt{firm\_nicring.Ring(size, refill)} avec \texttt{nic\_rx(seq, length)}, \texttt{poll(budget)},
\texttt{run(events)} et les compteurs \texttt{rx}, \texttt{drops}, \texttt{processed}, \texttt{doorbells}, \texttt{max\_owned},
\texttt{hash}~; \texttt{make\_events}. C++20 \texttt{firm::nicring::Ring} (\texttt{cpp/firm\_nicring.hpp}) et Rust
\texttt{firm\_nicring::Ring} avec les mêmes méthodes.
\bfield{Règles} La taille de l'anneau est une puissance de deux~; la carte n'écrase jamais un descripteur que l'hôte détient~; l'hôte traite dans l'ordre,
sans jamais dépasser son budget par scrutation, et rend les descripteurs par lots de \texttt{refill}, avec une sonnette par lot~; le hachage est FNV-1a calculé sur
les numéros de séquence et les longueurs des paquets traités, dans l'ordre de traitement.
\bfield{Tests d'acceptation} \texttt{code/firm/nicring/}~: le jeu de données partagé (20\,000 arrivées, quatre configurations) reproduit à l'identique en
Python, C++ et Rust~; un anneau plein perd des paquets, et un descripteur traité mais non rendu reste compté comme plein~; sous des entrelacements aléatoires,
chaque paquet accepté est traité une fois et dans l'ordre.
\bfield{Pour aller plus loin} Un anneau d'émission avec une file de complétion~; plusieurs anneaux avec une fonction d'aiguillage, pour reproduire le déséquilibre d'un hachage qui place
deux groupes actifs sur un même anneau.
\end{build}

\begin{omsources}
\item Documentation du noyau Linux~: ``NAPI'', ``Scaling in the Linux Networking Stack'', ``Timestamping'', ``AF\_XDP'';
\texttt{include/uapi/linux/ethtool.h} (\texttt{struct ethtool\_coalesce}); page de manuel \texttt{recvmmsg(2)}.
\item DPDK Programmer's Guide, ``Poll Mode Driver''; AMD, \emph{ef\_vi User Guide}; README d'OpenOnload.
\end{omsources}

\section{Exercices}

\begin{exercise}[$\star$]\label{exo:nw:network-cards-and-kernel-bypass:1}
Un anneau de 512 descripteurs, dont l'hôte en détient 40, reçoit 1.5 million de paquets par seconde pendant que son cœur est interrompu pendant
\qty{400}{\micro\second}. Calculer le nombre de paquets perdus.
\end{exercise}

\begin{exercise}[$\star$]\label{exo:nw:network-cards-and-kernel-bypass:2}
Avec un regroupement de 20 microsecondes, et des paquets à 200\,000 par seconde, calculer le nombre moyen de paquets servis par une interruption, et
le débit auquel arrivent les interruptions.
\end{exercise}

\begin{exercise}[$\star$]\label{exo:nw:network-cards-and-kernel-bypass:3}
D'après le \cref{tab:nw:network-cards-and-kernel-bypass:measured}, calculer quelle part de la médiane de la lecture bloquante est consommée après
l'entrée du datagramme dans le noyau.
\end{exercise}

\begin{exercise}[$\star\star$]\label{exo:nw:network-cards-and-kernel-bypass:4}
Expliquer pourquoi la répartition uniforme des flux par la \omterm{def:nw:network-cards-and-kernel-bypass:rss}{répartition côté réception} (RSS) joue contre un feed handler, et ce que change l'\omterm{def:nw:network-cards-and-kernel-bypass:rss}{aiguillage des flux} (flow steering).
\end{exercise}

\begin{exercise}[$\star\star$]\label{exo:nw:network-cards-and-kernel-bypass:5}
Avec les coûts du modèle, à partir de quel débit de paquets une interruption par paquet occupe-t-elle un cœur entier~? Et si chaque interruption servait 16 paquets~?
\end{exercise}

\begin{exercise}[$\star\star$]\label{exo:nw:network-cards-and-kernel-bypass:6}
Le busy polling a divisé la médiane par dix et a multiplié le 99\up{e}~centile par seize 
(\cref{tab:nw:network-cards-and-kernel-bypass:measured}). Expliquer ces deux effets, et dire ce qu'il faut changer sur un hôte de production.
\end{exercise}

\begin{exercise}[$\star\star\star$]\label{exo:nw:network-cards-and-kernel-bypass:7}
\emph{Programmation.} Avec \texttt{firm.nicring} et les événements de la fixture, trouver la plus petite taille d'anneau qui ne perd aucun paquet lorsque les descripteurs sont
rendus par lots de 32.
\end{exercise}

\begin{exercise}[$\star\star\star$]\label{exo:nw:network-cards-and-kernel-bypass:8}
\emph{Trouver la faille.} «~Nous avons déplacé la saisie d'ordres sur une interface brute en mode scrutation et écrit notre propre TCP minimal~: pas de minuteur de retransmission, puisque
les liaisons de colocation ne perdent pas de paquets. La latence a baissé de deux microsecondes.~»
\end{exercise}

\section{Problème~: Interruptions ou scrutation}

\begin{problem}[{Problème du week-end --- le prix d'un lecteur endormi}]\label{pb:nw:network-cards-and-kernel-bypass:1}
Un flux arrive à 200\,000 paquets par seconde sur un seul anneau de réception. Utiliser les coûts du modèle du chapitre~: \qty{2}{\micro\second} de CPU par
interruption, \qty{3}{\micro\second} de réveil, \qty{0.3}{\micro\second} par paquet, et une boucle de scrutation qui vérifie l'anneau toutes les
\qty{0.2}{\micro\second}.

\medskip\noindent\textbf{Partie I --- Une interruption par paquet.}
\begin{enumerate}
\item Quelle part d'un cœur occupent les interruptions et les paquets~?
\item Quelle est la plus petite latence qu'un paquet puisse avoir~?
\item La latence moyenne du modèle est de \qty{6.3}{\micro\second}. Pourquoi est-elle supérieure au minimum~?
\item À quel débit le cœur serait-il saturé~?
\end{enumerate}

\medskip\noindent\textbf{Partie II --- Un minuteur de 10 microsecondes.}
\begin{enumerate}[resume]
\item Combien de paquets une interruption sert-elle en moyenne~?
\item À quel débit les interruptions arrivent-elles~?
\item Quelle part du cœur est utilisée~?
\item La latence moyenne du modèle est de \qty{12.4}{\micro\second}. D'où vient la hausse~?
\end{enumerate}

\medskip\noindent\textbf{Partie III --- La scrutation.}
\begin{enumerate}[resume]
\item Quelle est l'attente moyenne avant la prochaine vérification de l'anneau~?
\item Quelle est la latence moyenne, et celle du modèle (\qty{0.41}{\micro\second})~?
\item Quelle part d'un cœur est utilisée~?
\item Que doit-il être vrai de ce cœur~?
\end{enumerate}

\medskip\noindent\textbf{Partie IV --- Le verdict.}
\begin{enumerate}[resume]
\item Énoncer le \emph{résultat nommé}~: la latence moyenne et le CPU du minuteur et de la scrutation à ce débit, ainsi que le débit auquel les seules interruptions occuperaient un cœur entier.
  Il faut donner ces grandeurs pour les deux cas.
\item Combien de microsecondes par paquet la scrutation économise-t-elle par rapport au minuteur~?
\item À quel débit de paquets le CPU du minuteur atteint-il un quart de cœur~?
\item Pourquoi les hôtes de trading scrutent-ils le chemin critique (hot path) et regroupent-ils le reste~?
\item Que faudrait-il mesurer sur la carte réelle avant de se fier au modèle~?
\item Que donne un \omterm{def:nw:network-cards-and-kernel-bypass:ring}{anneau de descripteurs} de 1\,024 entrées au cœur en scrutation, en temps, à ce débit~?
\item En une phrase~: que coûte un lecteur endormi~?
\item Laquelle des trois familles de contournement du chapitre utiliseriez-vous pour ce flux, et pourquoi~?
\end{enumerate}
\end{problem}

\section{Questions d'entretien}

\begin{interviewq}[$\star$ \iqroles{développeur}]\label{iq:nw:network-cards-and-kernel-bypass:1}
Décrivez ce qui se passe entre l'arrivée d'un paquet sur la carte et le retour de \texttt{recv}.
\end{interviewq}

\begin{interviewq}[$\star\star$ \iqroles{développeur}]\label{iq:nw:network-cards-and-kernel-bypass:2}
Qu'est-ce que le \omterm{def:nw:network-cards-and-kernel-bypass:coalesce}{regroupement des interruptions} (interrupt coalescing)~? Comment le configureriez-vous sur un hôte de données de marché~?
\end{interviewq}

\begin{interviewq}[$\star\star$ \iqroles{développeur}]\label{iq:nw:network-cards-and-kernel-bypass:3}
Qu'est-ce que le contournement du noyau (kernel bypass), quelles formes prend-il, et à quoi renonce-t-on~?
\end{interviewq}

\begin{interviewq}[$\star\star$ \iqroles{développeur}]\label{iq:nw:network-cards-and-kernel-bypass:4}
Votre feed handler en busy polling a une excellente médiane et un très mauvais 99.9\up{e}~centile. Que vérifiez-vous~?
\end{interviewq}

\begin{interviewq}[$\star\star\star$ \iqroles{développeur}]\label{iq:nw:network-cards-and-kernel-bypass:5}
Un \omterm{def:nw:network-cards-and-kernel-bypass:ring}{anneau de descripteurs} perd des paquets alors que la charge moyenne du gestionnaire est de 30~\%. Expliquez comment cela se produit et comment dimensionneriez-vous l'anneau.
\end{interviewq}

\begin{interviewq}[$\star\star$ \iqroles{développeur, chercheur}]\label{iq:nw:network-cards-and-kernel-bypass:6}
Pourquoi voudriez-vous des horodatages matériels de réception alors que vous horodatez déjà dans le noyau~?
\end{interviewq}
